\documentclass[10pt,a4paper]{amsart}
\usepackage[margin=19mm]{geometry}
\usepackage{amsmath,amssymb,mathtools}
\usepackage[scr=rsfs,cal=euler]{mathalfa}
\usepackage{xfrac}
\usepackage[unicode,hidelinks,bookmarks=false]{hyperref}
\newcommand{\ie}{i.e.\ }
\newcommand{\cf}{cf.\ }
\newcommand{\resp}{respectively\ }
\newcommand{\Subsection}{\subsection}
\providecommand{\nobreakdash}{\nobreak\hbox{-}}
\setlength{\emergencystretch}{2em}
\multlinegap0pt
\usepackage{amssymb}
\usepackage{xcolor}
\usepackage{cancel}
\newtheorem{corollary}{Corollary}
\newtheorem{proposition}{Proposition}
\newtheorem{claim}{Claim}
\newcommand{\maxidx}{\mathord{\xi}} % initial choice
\title{Optimal Transport and Generalized Lagrangian Mean Curvature Flows on Kim-McCann Metrics}
\author{Arunima Bhattacharya, Micah Warren,, Daniel Weser}
\date{Source: arXiv:2603.20498v1 (CC BY 4.0)}
\begin{document}
\maketitle

\noindent\emph{Source and licence.} Optimal Transport and Generalized Lagrangian Mean Curvature Flows on Kim-McCann Metrics. Version arXiv:2603.20498v1 (CC BY 4.0), distributed by arXiv under the Creative Commons Attribution 4.0 International licence, http://creativecommons.org/licenses/by/4.0/, as recorded in the arXiv licence metadata for this exact version and shown on the abstract page https://arxiv.org/abs/2603.20498v1. Full licence notice: \texttt{licenses/arxiv-2603.20498-CC-BY-4.0.txt}.\par\smallskip

\noindent\textit{Editorial scope.} Counted unit: the article's claim claim (source0.tex lines 3250--3260). The article's complete original proof of it is delivered with it (source0.tex lines 3262--3509, one \texttt{proof} environment of 248 lines). No mathematics is written, repaired or re-derived: the statement and the proof are the article's own lines, in the article's own notation, and no retained line is substituted or re-flowed.  Retained uncounted as the source-local context the proof needs: lines 2140--2142; lines 2778--2789; lines 3095--3109; lines 3511--3517; lines 3562--3569.  Nothing was omitted from any retained span, so the package carries no omission record.  No retained line contains a citation, so the wrapper carries no bibliography and no bibliography entry.  No retained line contains a figure environment, an image-inclusion command or drawing code, so no image is delivered and the package declares no asset dependency.  Editorial formatting only, in addition: an AMS \texttt{amsart} preamble that restates the article's own declarations and macros used by the retained lines, namely the environments \texttt{corollary}, \texttt{proposition}, \texttt{claim} on separate counters, together with the standard packages those lines need.  The retained environments print in this document as Corollary 1, Proposition 1, Claim 1 in order of appearance, whereas the article prints the same items with the numbering of its own full document.  The complete native source of the article as distributed in the arXiv e-print for this exact version is bundled unchanged as \texttt{sources/196850/source0.tex}.
\begin{corollary} \label{thetamaxprinciple}
    Given a generalized Lagrangian mean curvature flow, the value of $\theta$ is bounded above and below by its initial values. 
\end{corollary}

\begin{proposition}
\label{p1}
Let $S$ be a symmetric $(0,2)$-tensor on $M$. 
Suppose $p_0$ is a point where the function $R(V)$ achieves its maximum, 
and let $F_{\maxidx}$ be a maximizing vector. 
If $S$ is diagonal at $p_0$ with respect to normal coordinates 
$\{F_1,\ldots,F_n\}$ containing $F_{\maxidx}$, then
\[
\nabla_{F_i}\nabla_{F_i} S(F_{\maxidx},F_{\maxidx}) \le 0
\quad \text{for all } i .
\]
\end{proposition}

\begin{align}
\frac{d}{dt}S_{\maxidx\maxidx}
& - g^{mm}\nabla_{F_{m}}\nabla_{F_{m}}S(F_{\maxidx},F_{\maxidx})  \label{evolution S}\\
=& -n\hat{\nabla}_{\left( \nabla \psi \right) ^{\perp }}S(F_{\maxidx},F_{\maxidx})
   -2n\hat{S}\left( (\hat{\nabla}_{F_{\maxidx}}\left( \nabla \psi \right) ^{\perp })^{T},
     F_{\maxidx}\right)
   -2n\hat{S}\left( \nabla_{F_{\maxidx}}^{\perp }\left( \nabla \psi \right) ^{\perp },
     F_{\maxidx}\right) \nonumber \\
& - g^{mm}\hat{\nabla}_{F_{m}}\hat{\nabla}_{F_{m}}\hat{S}(F_{\maxidx},F_{\maxidx})
  + 2\hat{S}_{\maxidx\maxidx}\left( (\hat{\nabla}_{F_{\maxidx}}\vec{H})^{T}, F_{\maxidx}\right)
  - 4 g^{mm}\hat{\nabla}_{F_{m}}\hat{S}(A(F_{m},F_{\maxidx}),F_{\maxidx}) \nonumber \\
& - 2 g^{mm}\hat{S}(A(F_{m},F_{\maxidx}),A(F_{m},F_{\maxidx}))
  - 2 g^{mm}\sum \hat{R}(F_{m},F_{\maxidx},F_{m},KF_{l})\hat{S}(KF_{l},F_{\maxidx}) \nonumber \\
& + 2 g^{mm}A(F_{m},F_{\maxidx}) \cdot A(F_{m},F_{\maxidx}) \hat{S}_{\maxidx\maxidx}. \nonumber
\end{align}

\begin{claim}\label{a67}
\begin{equation*}
\hat{R}(F_{i},F_{\maxidx},F_{i},KF_{\maxidx})
= -\frac{\lambda_{\maxidx}}{\lambda_{i}} \hat{R}_{i\overline{\maxidx} i\overline{\maxidx}}
  + \frac{\lambda_{i}}{\lambda_{\maxidx}} \hat{R}_{\overline{i} \maxidx \overline{i} \maxidx}.
\end{equation*}
\end{claim}

\begin{claim} \label{annoying}%
\[
2\hat{S}\left(  \nabla_{F_{\maxidx}}^{\perp}n\left(  \nabla\psi\right)  ^{\perp}
,F_{\maxidx}\right)  \leq\left\Vert F_{\maxidx}\right\Vert _{S}\left\Vert \left(
\hat{\nabla}_{F_{\maxidx}}n\left(  \nabla\psi\right)  ^{\perp}\right)  ^{\perp}
\right\Vert _{\hat{S}}.
\]
\end{claim}

\begin{claim}
There exist constants $\bar C>0$ and $\kappa>0$ such that the following holds. 
If $S_{\maxidx\maxidx}\ge 10$ and $F_{\maxidx}$ is the maximizer of the ratio $\mathcal{R}(F_{\maxidx})$, then at the maximum point we have
\[
\partial_t Q(F_{\maxidx},F_{\maxidx})
\le 
\bar C\, S_{\maxidx\maxidx}^2
-
\kappa\, S_{\maxidx\maxidx}^{2+\frac{1}{n-1}}.
\]
\end{claim}

\begin{proof}
We start by computing
\[
\partial_{t}Q(F_{\maxidx},F_{\maxidx})=\partial_{t}\left(  \frac{S(F_{\maxidx},F_{\maxidx},t)}%
{g(F_{\maxidx},F_{\maxidx},t)}\right)  \Big|_{x_{0},t_{0}}=S_{\maxidx\maxidx,t}-S_{\maxidx\maxidx}g_{\maxidx\maxidx,t}.
\]
We subtract the {nonpositive} quantity (nonpositive by Proposition \ref{p1}) to obtain
\begin{align*}
\partial_{t}Q(F_{\maxidx},F_{\maxidx})  & \leq\partial_{t}Q(F_{\maxidx},F_{\maxidx})-\sum_{i}%
\nabla_{F_{i}}\nabla_{F_{i}}S(F_{\maxidx},F_{\maxidx})\\
&  =S_{\maxidx\maxidx,t}-\sum_{i}\nabla_{F_{i}}\nabla_{F_{i}}S(F_{\maxidx},F_{\maxidx})-S_{\maxidx\maxidx}%
g_{\maxidx\maxidx,t}.
\end{align*}

Next, we plug in \eqref{evolution S}:
\begin{align*}
\partial_{t}Q(F_{\maxidx},F_{\maxidx})
& \leq -\hat{\nabla}_{n\left( \nabla\psi\right)^{\perp}}
      S(F_{\maxidx},F_{\maxidx})
      -2\hat{S}\left( \nabla_{F_{\maxidx}}^{\perp} n\left(\nabla\psi\right)^{\perp},
      F_{\maxidx}\right)
      -\sum_{i}\hat{\nabla}_{F_{i}}\hat{\nabla}_{F_{i}}
      \hat{S}(F_{\maxidx},F_{\maxidx})\\
& \quad -4\sum_{i}\hat{\nabla}_{F_{i}}
      \hat{S}(A(F_{i},F_{\maxidx}),F_{\maxidx})\\
& \quad +2\hat{S}_{\maxidx\maxidx}\left(
        (\hat{\nabla}_{F_{\maxidx}}\vec{H})^{T}\cdot F_{\maxidx}\right)
      -2\hat{S}\left(
        (\hat{\nabla}_{F_{\maxidx}}n\left( \nabla\psi\right)^{\perp})^{T},
        F_{\maxidx}\right)
      -S_{\maxidx\maxidx} g_{\maxidx\maxidx,t}\\
& \quad -2\sum_{i}\hat{S}(A(F_{i},F_{\maxidx}),A(F_{i},F_{\maxidx}))
      -2\sum_{i,l}\hat{R}(F_{i},F_{\maxidx},F_{i},KF_{l})
         \hat{S}(KF_{l},F_{\maxidx})\\
& \quad +2\sum_{i}
      A(F_{i},F_{\maxidx})\cdot A(F_{i},F_{\maxidx})\hat{S}_{\maxidx\maxidx}.
\end{align*}
We rearrange the above quantities as follows:\begin{align*}
&\quad +2\sum_{i}A(F_{i},F_{\maxidx})\cdot A(F_{i},F_{\maxidx})\hat{S}_{\maxidx\maxidx}
\tag{1}\label{eq:L1}\\
&\quad +2\hat{S}_{\maxidx\maxidx}\left( (\hat{\nabla}_{F_{\maxidx}}\vec{H})^{T}\cdot F_{\maxidx}\right)
   -2\hat{S}\left( (\hat{\nabla}_{F_{\maxidx}}n(\nabla\psi)^{\perp})^{T},F_{\maxidx}\right)
   -S_{\maxidx\maxidx}g_{\maxidx\maxidx,t}
\tag{2}\label{eq:L2}\\
&\quad -4\sum_{i}\hat{\nabla}_{F_{i}}\hat{S}(A(F_{i},F_{\maxidx}),F_{\maxidx})
   \;-\;2\sum_{i}\hat{S}(A(F_{i},F_{\maxidx}),A(F_{i},F_{\maxidx}))
\tag{3}\label{eq:L3}\\
&\quad -\hat{\nabla}_{n(\nabla\psi)^{\perp}}S(F_{\maxidx},F_{\maxidx})
   -\sum_{i}\hat{\nabla}_{F_{i}}\hat{\nabla}_{F_{i}}\hat{S}(F_{\maxidx},F_{\maxidx})
\tag{4}\label{eq:L4}\\
&\quad -2\hat{S}\left( \nabla_{F_{\maxidx}}^{\perp}n(\nabla\psi)^{\perp},F_{\maxidx}\right)
\tag{5}\label{eq:L5}\\
&\quad -2\sum_{i,l}\hat{R}(F_{i},F_{\maxidx},F_{i},KF_{l})\hat{S}(KF_{l},F_{\maxidx})
\tag{6}\label{eq:L6}
\end{align*}

We will deal with these lines in order. \ For lines 1-5 we will show that each
expression is bounded by%
\[
C(1+S_{\maxidx\maxidx}^{2})
\]
for some controlled constant $C.$

\textit{Line 1:}\ \ \ Note that this is the sum of metric pairings of two time-like
vectors, so this is non positive.

\textit{Line 2:} Note that
\begin{align*}
\frac{d}{dt}g_{\maxidx\maxidx} 
&=2h(\hat{\nabla}_{F_{\maxidx}}\left(  \vec{H}-n\left(
\nabla\psi\right)  ^{\perp}\right)  ,F_{\maxidx})\\
&=2(\hat{\nabla}_{F_{\maxidx}}\vec{H})^{T}\cdot F_{\maxidx}
   -2\left(  \hat{\nabla}%
_{F_{\maxidx}}n\left(  \nabla\psi\right)  ^{\perp}\right)  ^{T}\cdot F_{\maxidx}
\end{align*}
and that
\[
\hat{S}\left(  (\hat{\nabla}_{F_{\maxidx}}n\left(  \nabla\psi\right)  ^{\perp}%
)^{T},F_{\maxidx}\right)
=(\hat{\nabla}_{F_{\maxidx}}n\left(  \nabla\psi\right)  ^{\perp})^{T}
\cdot F_{\maxidx}\, S(F_{\maxidx},F_{\maxidx}).
\]

So
\begin{align*}
&\quad 
+2\hat{S}_{\maxidx\maxidx}\left( (\hat{\nabla}_{F_{\maxidx}}\vec{H})^{T}\!\cdot F_{\maxidx}\right)
-2\hat{S}\!\left( (\hat{\nabla}_{F_{\maxidx}} n(\nabla\psi)^{\perp})^{T},\,F_{\maxidx}\right)
-S_{\maxidx\maxidx} g_{\maxidx\maxidx,t}
\\[0.5em]
&= 
2\hat{S}_{\maxidx\maxidx}\left( (\hat{\nabla}_{F_{\maxidx}}\vec{H})^{T}\!\cdot F_{\maxidx}\right)
-2\hat{S}\!\left( (\hat{\nabla}_{F_{\maxidx}} n(\nabla\psi)^{\perp})^{T},\,F_{\maxidx}\right)
\\
&\quad 
-S_{\maxidx\maxidx} 2(\hat{\nabla}_{F_{\maxidx}}\vec{H})^{T}\!\cdot F_{\maxidx}
+S_{\maxidx\maxidx} 2\left( \hat{\nabla}_{F_{\maxidx}} n(\nabla\psi)^{\perp} \right)^{\perp}\!\cdot F_{\maxidx}
\\[0.5em]
&=0.
\end{align*}

\textit{Line 3:} Note that as a fixed tensor on the ambient manifold, we have
\[
\hat{\nabla}_{F_{i}}\hat{S}(X,Y)
\leq C_{1}\left\Vert F_{i}\right\Vert
_{\hat{S}}\left\Vert X\right\Vert _{\hat{S}}\left\Vert Y\right\Vert _{\hat{S}},
\]
so
\[
-4\sum_{i}\hat{\nabla}_{F_{i}}\hat{S}(A(F_{i},F_{\maxidx}),F_{\maxidx})
\leq
2\sum_{i}\hat{S}(A(F_{i},F_{\maxidx}),A(F_{i},F_{\maxidx}))
+2\sum_{i}C_{1}^{2}\hat{S}(F_{i},F_{i})\hat{S}(F_{\maxidx},F_{\maxidx})
\]

hence
\begin{align*}
-4\sum_{i}\hat{\nabla}_{F_{i}}\hat{S}(A(F_{i},F_{\maxidx}),F_{\maxidx}
-2\sum_{i}\hat{S}(A(F_{i},F_{\maxidx}),A(F_{i},F_{\maxidx}))  
&\leq 2\sum_{i}C_{1}^{2}\hat{S}(F_{i},F_{i})\hat{S}(F_{\maxidx},F_{\maxidx})\\
&\leq C S_{\maxidx\maxidx}^{2}.
\end{align*}

\textit{Line 4:}
\[
-\hat{\nabla}_{n\left(  \nabla\psi\right)  ^{\perp}}
S(F_{\maxidx},F_{\maxidx})
-\sum_{i}\hat{\nabla}_{F_{i}}\hat{\nabla}_{F_{i}}\hat{S}(F_{\maxidx},F_{\maxidx}).
\]
Note that the first term involves a projection, so
\[
\left\Vert n\left(  \nabla\psi\right)  ^{\perp}\right\Vert _{\hat{S}}
=\left\Vert
h(n\left(  \nabla\psi\right)  ^{\perp},KF_{l})g^{lm}KF_{m}\right\Vert_{\hat{S}} 
\leq S_{\maxidx\maxidx}.
\]
This shows that the first term is bounded by $S_{\maxidx\maxidx}^{2}.$ 
The second is an ambient 4 tensor, so is bounded by$S_{\maxidx\maxidx}^{2}.$\\

\textit{Line 5:} This is Claim \ref{annoying} below.

We finally claim that the contribution from line 6 is bounded above by
\[
-\kappa S_{\maxidx\maxidx}^{2+\frac{1}{n-1}} + C S_{\maxidx\maxidx}^{2},
\]
for some constant $\kappa>0$, bounded away from zero, 
depending on the cross-curvature condition.

To set this up, recall
\[
m+\bar{m}=\delta_{ij}+\delta_{\bar{\imath}
\bar{j}}
\]
at the point, and the orthonormal basis with respect to $g$ has the
expression
\[
F_i=\frac{1}{\sqrt{\lambda_{i}}}E_{i}
  +\sqrt{\lambda_{i}}E_{\bar{\imath}}.
\]
Next, we observe the term
\[
\sum_{i,l}\hat{R}(F_{i},F_{\maxidx},F_{i},KF_{l})\hat{S}(KF_{l},F_{\maxidx})
\]
and note that
\[
\hat{S}(KF_{l},F_{\maxidx})
=\left(  \frac{1}{\lambda_{\maxidx}}-\lambda_{\maxidx}\right)\delta_{l\maxidx}.
\]
and 
\[
\hat{R}(F_{i},F_{\maxidx},F_{i},KF_{\maxidx})
=-\frac{\lambda_{\maxidx}}{\lambda_{i}}\hat{R}_{i\overline{\maxidx} i\overline{\maxidx}}
+\frac{\lambda_{i}}{\lambda_{\maxidx}}\hat{R}_{\overline{\imath}\maxidx\overline{\imath}\maxidx}.
\]

Thus (see Claim \ref{a67} below) 
\begin{align*}
\sum_{i,l}\hat{R}(F_{i},F_{\maxidx},F_{i},KF_{l})\hat{S}(KF_{l},F_{\maxidx})
&=\sum_{i}\left(
  -\frac{\lambda_{\maxidx}}{\lambda_{i}}\hat{R}_{i\overline{\maxidx}i\overline{\maxidx}}
  +\frac{\lambda_{i}}{\lambda_{\maxidx}}\hat{R}_{\overline{\imath}\maxidx\overline{\imath}\maxidx}
\right)
\left(  \frac{1}{\lambda_{\maxidx}}-\lambda_{\maxidx}\right)  \\
&=-\sum_{i}\frac{1}{\lambda_{i}}\hat{R}_{i\overline{\maxidx}i\overline{\maxidx}}
  +\lambda_{\maxidx}^{2}\sum_{i}\frac{1}{\lambda_{i}}\hat{R}_{i\overline{\maxidx}i\overline{\maxidx}}\\
&\qquad+\frac{1}{\lambda_{\maxidx}^{2}}\sum_{i}\lambda_{i}\hat{R}_{\overline{\imath}%
\maxidx\overline{\imath}\maxidx}-\sum_{i}\lambda_{i}\hat{R}_{\overline{\imath}\maxidx\overline{\imath}\maxidx}.
\end{align*}

If $\maxidx=1$ then the term
\[
\lambda_{1}^{2} \sum_{i}\frac{1}{\lambda_{i}}\hat{R}_{i\bar{1}i\bar{1}}
\]
dominates. Recall (here we use Corollary \ref{thetamaxprinciple}) 
\[
\det DT=\lambda_{1}\cdots\lambda_{n}
=\frac{e^{-2\theta}\rho}{\bar{\rho}}
\leq\Lambda_{0}
\]
so
\[
\Lambda_{0}\geq\lambda_{1}(\lambda_{n})^{\,n-1}
\]
and hence
\[
\frac{1}{\lambda_{n}}
\geq\left(  \frac{\lambda_{1}}{\Lambda_{0}}\right)^{\!\frac{1}{n-1}}.
\]

Thus
\[
\lambda_{1}^{2}\sum_{i}\frac{1}{\lambda_{i}}\hat{R}_{i\bar{1}i\bar{1}}
\geq
\lambda_{1}^{2}\left(  \frac{\lambda_{1}}{\Lambda_{0}}\right)^{\!\frac{1}{n-1}}
\hat{R}_{n\bar{1}n\bar{1}}
\geq
\kappa S_{\maxidx\maxidx}^{2+\frac{1}{n-1}}-C.
\]

Of the remaining terms,
\[
\frac{1}{\lambda_{\maxidx}^{2}}\sum_{i}\lambda_{i}\hat{R}_{\overline{\imath}\maxidx\overline{\imath}\maxidx}
\]
is positive, and the other two are bounded by $\sum_{i}S_{ii}.$

If $\maxidx=n,$ then
\[
\det DT=\lambda_{1}\cdots\lambda_{n}
=\frac{e^{-2\theta}\rho}{\bar{\rho}}
\geq\frac{1}{\Lambda_{0}}
\]
so
\[
\lambda_{1}\geq
\frac{1}{\Lambda_{0}(\lambda_{n})^{\,n-1}}
\]
and
\begin{align*}
\frac{1}{\lambda_{n}^{2}}\sum_{i}\lambda_{i}\hat{R}_{\bar{\imath}n\bar{\imath}n}
&\geq
\frac{1}{\lambda_{n}^{2}}
\frac{1}{\Lambda_{0}(\lambda_{n})^{\,n-1}}
\hat{R}_{\bar{1}n\bar{1}n}\\
&\geq
\kappa S_{\maxidx\maxidx}^{2+\frac{1}{n-1}}-C.
\end{align*}

\end{proof}

\end{document}
